% %%%%%%%%%%%%%%%%%%%%%%%%%%%%%%%%%%%%%%%%%%%%%%%%%%%%%%%%%%%%
% CHAPTER 5 - CONCLUSION AND RECOMMENDATIONS (guideline 2.5) ~9 pages
% Generated from Chapter5.src.tex: review-ID citations (\pc, \tc, \pcs)
% are converted to biblatex citations with verified first pages.
% %%%%%%%%%%%%%%%%%%%%%%%%%%%%%%%%%%%%%%%%%%%%%%%%%%%%%%%%%%%%

\chapter{Conclusion and Recommendations}
\label{ch:conclusion}

This final chapter brings the thesis to a close. It first summarises the study and answers the three research questions (Section~\ref{sec:summary}). It then outlines the theoretical contributions (Section~\ref{sec:contributions}) and derives recommendations for managers of e-commerce and platform-based businesses (Section~\ref{sec:recommendations}). Section~\ref{sec:future} reflects critically on the limitations of the thesis and identifies avenues for future research, before Section~\ref{sec:final} concludes.

% ------------------------------------------------------------
\section{Summary and Answers to Research Questions}
\label{sec:summary}

The starting point of this thesis was a tension observed in Chapter~\ref{ch:intro}. On the one hand, AI is widely regarded as a powerful source of value in data-rich digital businesses. On the other hand, AI technologies are becoming broadly accessible, and many firms struggle to turn their AI investments into measurable returns. The aim of the thesis was therefore to examine, on the basis of the existing academic literature, how AI capabilities contribute to sustainable competitive advantage in e-commerce and platform-based digital businesses, and which role dynamic capabilities play in this process.

To achieve this aim, the thesis combined the resource-based view and dynamic capabilities theory into a conceptual framework (Chapter~\ref{ch:literature}) and conducted a systematic literature review (Chapter~\ref{ch:methodology}). Two searches in the EBSCO Discovery Service identified 612 records, of which 109 studies were included after de-duplication, screening and eligibility assessment. Eighteen of these are core studies from e-commerce, online retail and platform contexts, and 91 are supporting studies from other industries. The evidence was extracted with a structured framework, checked through an independent second coding and, where accessible, against the full texts, assessed for quality and synthesised thematically along the research questions (Chapter~\ref{ch:results}). Table~\ref{tab:rq-answers} summarises the answers to the three research questions.

\begin{table}[H]
  \centering
  \caption{Summary of Answers to the Research Questions}
  \label{tab:rq-answers}
  \small\setstretch{1.0}\sloppy
  \begin{tabularx}{\textwidth}{>{\raggedright\arraybackslash}p{3.4cm}>{\raggedright\arraybackslash}X>{\raggedright\arraybackslash}p{1.8cm}}
    \toprule
    \textbf{Research question} & \textbf{Answer} & \textbf{Section} \\
    \midrule
    RQ1: How are AI capabilities associated with competitive advantage? & Positively, in the large majority of studies and in the core e-commerce studies, but mostly indirectly; direct effects often become weak once mediators are considered. Evidence concerns mainly performance and competitive advantage, rarely sustainable advantage. & \ref{sec:results-ai-ca}, \ref{sec:proposition-evaluation} \\
    \addlinespace
    RQ2: How do sensing, seizing and transforming relate to this relationship? & Dynamic capabilities are the main mechanism through which AI capabilities create value. Seizing (agility, innovation) dominates the evidence, transforming and sensing are studied much less often, although transforming appears most relevant for sustainability. & \ref{sec:results-ai-dc}, \ref{sec:mediation} \\
    \addlinespace
    RQ3: Which factors enable or constrain sustainable advantage through AI? & Integration with processes and complementary assets, environmental dynamism, data and infrastructure, managerial cognition, trust, ethics and regulation, firm resources, platform conditions and imitability. Their effect on the persistence of advantage over time is not yet established. & \ref{sec:proposition-evaluation}, \ref{sec:discussion} \\
    \bottomrule
  \end{tabularx}
  \source{Author's own work}
\end{table}

With regard to \emph{RQ1}, the review shows that AI capabilities are associated with competitive advantage and firm performance in a consistent way. Sixty-three of the 79 studies that examine this link report a positive association, and the same holds for 12 of the 14 core studies from e-commerce and platform contexts, while the other two describe the link as conditional. However, the association is predominantly indirect. In several studies, the direct effect of AI capabilities is not significant once mediating capabilities are taken into account \pcs{R003,R011,R018,R120,R059}. Furthermore, most studies measure firm performance or competitive advantage at one point in time. Firm performance and competitive advantage were coded separately from sustainable competitive advantage, and only 11 studies examine the latter; none of them is set in e-commerce. The answer to RQ1 is therefore that AI capabilities are associated with competitive advantage, but that the existing evidence does not yet demonstrate that this advantage is sustainable.

With regard to \emph{RQ2}, the review provides strong support for the dynamic capabilities perspective. Seventy-four of the 79 studies that address the link report that AI capabilities strengthen sensing, seizing or transforming, or that these capabilities explain how AI creates value. The mechanism is, however, uneven. Seizing, in particular in the form of organisational, strategic and supply chain agility and of service and marketing innovation, is addressed most often. Transforming, in the form of business model innovation and reconfiguration, is addressed less often, but the evidence suggests that it is decisive for moving beyond efficiency gains towards new value propositions \pc{O11}. Sensing is addressed about as often as transforming and far less often than seizing, although AI appears naturally suited to support it. The answer to RQ2 is therefore that dynamic capabilities, and in particular transforming capabilities, are the mechanism through which AI capabilities can be converted into a sustainable advantage.

With regard to \emph{RQ3}, the review identifies eight families of factors that enable or constrain the creation of advantage through AI capabilities. The most frequently reported are the integration of AI into business processes and the availability of complementary assets, followed by environmental dynamism, data quality and digital infrastructure, managerial cognition, and trust, ethics and regulation. Firm resources and characteristics, platform and ecosystem conditions and the imitability of AI applications complete the picture. These factors explain why identical AI investments lead to different results in different firms. Because the evidence base contains few longitudinal studies, however, it remains open which of these factors determine whether an AI-based advantage persists over time.

Taken together, the central conclusion of the thesis is that AI capabilities are an important, but not a sufficient, condition for sustainable competitive advantage. AI capabilities create value when firms convert them into dynamic capabilities, especially into the ability to transform their business models and resource base, and when complementary assets, governance and managerial cognition support this conversion. In e-commerce and platform businesses, where AI tools are widely accessible and competitors can quickly imitate efficiency gains, this conversion is likely to be the main source of a lasting advantage.

% ------------------------------------------------------------
\section{Theoretical Contributions}
\label{sec:contributions}

The thesis makes five contributions to research at the intersection of strategic management and information systems.

First, it integrates the resource-based view and dynamic capabilities theory into a coherent explanation of how AI creates competitive advantage. The review shows that AI capabilities rarely meet the criteria of rareness and inimitability that \textcite[99]{barney1991} identified as conditions for sustained competitive advantage, because AI tools are becoming democratised \pc{O16}. Their value lies instead in the way they are situated in firm-specific activities \pc{R087} and in the dynamic capabilities through which they are deployed and renewed. This supports the argument of \textcite[1319]{teece2007} that sustainable advantage rests on the capabilities that orchestrate and reconfigure resources, and it applies this argument specifically to AI.

Second, the thesis contributes to conceptual clarity by distinguishing systematically between firm performance, competitive advantage and sustainable competitive advantage. By coding these outcomes separately, the review reveals that the literature frequently treats performance gains as evidence of competitive advantage, while sustainable advantage is rarely measured. This distinction, which responds to the conceptual ambiguity noted by \textcite[1]{wu2026}, is a precondition for cumulative research on AI and long-term competitiveness.

Third, the thesis refines the dynamic capabilities mechanism for AI. Rather than treating dynamic capabilities as an aggregate construct, it shows that sensing, seizing and transforming are addressed very unevenly and that their contributions differ. The finding that seizing dominates the evidence, while transforming appears most relevant for sustainability, connects the AI literature with the argument of \textcite[1105]{eisenhardt2000} that dynamic capabilities often show commonalities across firms. If many firms use AI to improve similar seizing routines, the resulting gains are likely to converge, and only the reconfiguration of business models and resource bases can create lasting differences.

Fourth, the thesis provides a structured map of the contingencies that shape the effect of AI capabilities. Previous studies have examined individual moderators in isolation. The review integrates them into eight families and adds two groups that are particularly relevant for digital business but have received little attention in the capability literature: trust, ethics and regulation, and platform and ecosystem conditions such as data network effects and algorithm competition \pcs{O15,O14}.

Fifth, the evidence-informed conceptual model in Figure~\ref{fig:final-model} offers a basis for future empirical research. It specifies the expected pathways from AI capabilities through the three dynamic capabilities to competitive and sustainable competitive advantage, and it identifies the boundary conditions that should be considered. In addition, the distinction between core and supporting evidence, combined with a transparent quality assessment, provides a methodological template for reviews in fields where context-specific evidence is still scarce.

% ------------------------------------------------------------
\section{Managerial Recommendations}
\label{sec:recommendations}

The findings lead to six recommendations for managers of e-commerce and platform-based businesses who want to turn AI into a lasting competitive advantage. They are summarised in Table~\ref{tab:recommendations} and explained below.

\begin{table}[H]
  \centering
  \caption{Managerial Recommendations}
  \label{tab:recommendations}
  \small\setstretch{1.0}\sloppy
  \begin{tabularx}{\textwidth}{>{\raggedright\arraybackslash}p{4.2cm}>{\raggedright\arraybackslash}X>{\raggedright\arraybackslash}p{3.4cm}}
    \toprule
    \textbf{Recommendation} & \textbf{Rationale} & \textbf{Supporting evidence} \\
    \midrule
    1. Treat AI as an organisational capability, not as a tool & AI tools are accessible to competitors; value lies in firm-specific combinations & \pcs{O16,R087,O01} \\
    \addlinespace
    2. Invest in complementary assets & Data quality, infrastructure and integration determine the effect of AI & \pcs{R050,R111,R083} \\
    \addlinespace
    3. Use AI to transform, not only to optimise & Efficiency gains are easily imitated; new value propositions last longer & \pcs{O11,O09,R037} \\
    \addlinespace
    4. Build agility and resilience & Agility and resilience are the main channels from AI to performance & \pcs{O04,R018,R108} \\
    \addlinespace
    5. Make trust and governance part of the AI strategy & Customer data and trust underpin AI-based advantage in e-commerce & \pcs{R148,R133,R043} \\
    \addlinespace
    6. Develop managerial AI cognition and monitor competition & Managers' framing of AI and competitors' algorithms shape the persistence of advantage & \pcs{O07,O14,R129} \\
    \bottomrule
  \end{tabularx}
  \source{Author's own work based on the findings in Chapter~\ref{ch:results}}
\end{table}

\emph{Treat AI as an organisational capability.} Managers should avoid equating the purchase or adoption of AI tools with an AI strategy. Since many AI applications, including generative AI, are available to all competitors, they are unlikely to create a lasting difference on their own. Instead, managers should build an AI capability in which several AI-specific resources are combined in a firm-specific way \pc{O01} and situate AI in the specific activities and decisions of their business \pc{R087}.

\emph{Invest in complementary assets.} The effect of AI capabilities depends on the quality of data \pc{R050}, on digital infrastructure \pc{R111} and on the integration of AI with suppliers and customers \pc{R083}. For e-commerce firms, this means that investments in AI should be accompanied by investments in data governance, system integration and logistics, which are often less visible but decisive for success.

\emph{Use AI to transform, not only to optimise.} Many e-commerce firms use AI mainly to increase efficiency, for example in logistics, pricing or customer service. Such gains are valuable but can be imitated quickly. The evidence suggests that firms should also use AI to develop new value propositions and business models, for example AI-based services, personalised offerings or platform-based ecosystems, because reconfiguration is required to escape the efficiency trap \pc{O11} and to scale AI \pc{O09}.

\emph{Build agility and resilience.} Agility and resilience are among the most frequent channels through which AI capabilities improve performance and competitive advantage \pcs{O04,R018}. Managers should therefore use AI-generated insights to accelerate decision making and to adapt supply chains, assortments and marketing quickly to changing conditions, and they should design organisational structures and routines that allow such adaptation.

\emph{Make trust and governance part of the AI strategy.} Customer-facing AI in e-commerce depends on customers' willingness to share data. Managers should therefore communicate transparently about their use of AI, address privacy, bias and explainability concerns \pc{R133}, implement ethical AI governance \pc{R043} and avoid overstating their AI capabilities, since AI washing can damage trust and reputation \pc{R148}.

\emph{Develop managerial AI cognition and monitor competition.} The way in which managers interpret AI, as a tool, an actor or a strategic resource, shapes whether a firm achieves fleeting efficiency gains or sustained advantages \pc{O07}. Managers should therefore invest in their own AI literacy and in that of their teams. At the same time, they should monitor competitors' AI capabilities closely, since an entrant with a superior algorithm can overcome even a large user base \pc{O14}, and the performance effects of AI depend on how quickly competitors adopt similar capabilities \pc{R129}.

The six recommendations are not independent measures but can be understood as a sequence. In the short term, firms should establish the foundations of an AI capability, namely reliable data, digital infrastructure and the integration of AI into core processes, and they should define governance rules for the responsible use of customer data. In the medium term, they should use these foundations to build agility, for example by accelerating decisions on assortments, prices and marketing campaigns, and by using AI to develop new services and customer experiences. In the long term, they should use the resulting insights and capabilities to transform their business models and ecosystem relationships, since this is where the evidence locates the most durable sources of advantage. Throughout this sequence, managers should regularly assess whether their AI-based benefits are still distinctive or whether competitors have caught up, because this assessment indicates when the next cycle of renewal is required.

% ------------------------------------------------------------
\section{Limitations and Future Research}
\label{sec:future}

The findings of this thesis are subject to several limitations, some of which arise from the review method and some from the state of the evidence. Regarding the method, the search was conducted in a single discovery service, restricted to English-language, peer-reviewed journal articles, and the screening was carried out by a single reviewer. Most importantly, the data extraction for 80 of the 109 studies was based on the abstracts, complemented by bibliographic information; only the 29 accessible open-access studies could be checked against their full texts. Abstracts report the main findings but omit many details, for example on measurement, sample characteristics or limitations. Whether limitations are discussed could therefore be assessed for only 36 studies and was reported separately from the quality score instead of being scored as not met, and some findings may be simplified. These limitations were mitigated by explicit criteria, reason codes, a documented screening log and an AI-assisted second coding with adjudication of all disagreements, but they should be addressed in future work through a complete full-text analysis and independent coding by several human reviewers.

Regarding the evidence, the review is limited by the characteristics of the primary studies. Most of them are cross-sectional surveys based on managers' perceptions, which do not allow causal conclusions and may be affected by common method bias. The constructs used to measure AI capabilities, dynamic capabilities and competitive advantage differ considerably, which prevents a quantitative comparison of effect sizes. The evidence is geographically concentrated in China, the Middle East and South Asia, and none of the studies that report their context is set in Germany. Finally, three quarters of the included studies were published in 2025 and 2026, so that the evidence base is likely to change quickly.

Looking back on the research process, several decisions of the author also shaped the results and deserve critical reflection. The two-step search strategy kept the review manageable, but it relied on the term \emph{AI capability} in the primary search, which is why 16 relevant studies had to be added through other methods; a broader string with a larger review team might have found further studies. The decision not to use journal rankings as an inclusion criterion increased the coverage of the evidence, but it also brought studies of very different quality into the sample, which made the quality assessment more important and the synthesis more complex. The distinction between core and supporting studies proved useful for separating context-specific from general evidence, but the boundary between the two tiers is not always sharp, for example for retail or supply chain studies with an e-commerce component. Finally, the conceptual framework derived in Chapter~\ref{ch:literature} guided the coding and may have directed attention towards dynamic capabilities at the expense of alternative explanations, such as institutional or behavioural theories, which some of the included studies also use. Being aware of these choices is important for interpreting the findings, and future reviews could address them by combining several search strategies, coders and theoretical lenses.

These limitations point to several avenues for future research:
\begin{itemize}
  \item \emph{Longitudinal research on sustainability.} Future studies should observe firms over several years to examine whether and how long AI-based advantages persist, using panel data or repeated surveys, and should measure sustainable competitive advantage directly rather than inferring it from performance.
  \item \emph{E-commerce and platform-specific studies.} Since only 18 of the 109 studies are set in e-commerce or platform contexts and none examines sustainable advantage, research should test the evidence-informed model in these contexts, including the role of data network effects and algorithm competition.
  \item \emph{European and German evidence.} Studies of European and particularly German e-commerce firms would show whether the findings hold in markets with different regulatory conditions, such as stricter data protection.
  \item \emph{Differentiated dynamic capabilities.} Future research should examine sensing, seizing and transforming separately and compare their contributions, with particular attention to the under-researched sensing capability and to the transforming capabilities that appear most relevant for sustainability.
  \item \emph{Generative and agentic AI.} As generative and agentic AI diffuse rapidly, research should examine whether these technologies change the conditions of imitability and the sources of advantage identified in this review.
  \item \emph{Meta-analysis and full-text reviews.} Once the evidence base matures, a meta-analysis could quantify the strength of the relationships, and a full-text review with several coders could validate the findings of this abstract-based synthesis.
\end{itemize}

% ------------------------------------------------------------
\section{Final Conclusion}
\label{sec:final}

This thesis set out to understand how AI capabilities contribute to sustainable competitive advantage in e-commerce and platform-based digital business. Based on a systematic review of 109 studies, it concludes that AI capabilities are consistently associated with competitive advantage, but that this advantage arises mainly indirectly, through dynamic capabilities such as agility, innovation and the transformation of business models, and depends on complementary assets, governance, managerial cognition and the competitive environment. At the same time, the review shows that the literature has so far focused on performance gains at a single point in time and has rarely examined whether AI-based advantages last, especially in e-commerce.

For research, the thesis offers an evidence-informed model that links AI capabilities, sensing, seizing and transforming, and sustainable competitive advantage, and identifies the conditions under which this link holds. For practice, it suggests that e-commerce and platform firms should not ask whether to use AI, since AI is quickly becoming a standard tool, but how to build the capabilities that allow them to renew their advantage faster than competitors can imitate it. In a market in which algorithms, data and customer expectations change continuously, the most durable advantage may not be a particular AI capability, but the dynamic capability to keep reinventing how AI creates value.
